% p1-13-limitations

\section{P1 §13. Limitations and Scope of Validity}

13. Limitations and Scope of Validity

  Figure (fig-p1-ch13-limitations). Limitations triptych --- Grammar non-uniqueness / Finite dataset
                                       / Asymptotic saturation.

  13.1 Grammar Dependence and Non-Uniqueness
  All compressibility properties documented here are conditional on the specific choice of
  grammar G$\varphi$ = langle $\varphi$, pi, e, 3, Z rangle. Different bases induce inequivalent
  hypothesis spaces, and a finding of compressibility under G$\varphi$ does not imply that {$\varphi$, pi,
  e, 3} is uniquely suited to the task. Future work should perform a systematic basis
  comparison using the same MDL framework.

  13.2 Residual Selection and Search-Induced Bias
  The exhaustive search over S(C) generates a minimum residual dC(X) that is an
  extreme-value statistic. Despite the null-model controls, residual selection bias may
  persist due to:

  - Anisotropy of the log-lattice induced by the specific values of log 3, log $\varphi$, log pi, log e.
  - Clustering of grammar values near dense regions that accidentally overlap with the
  empirical distribution of physical constants.
  - Correlations between residuals across different constants in the finite, non-random
  target set T.

  The Mperm null model partially addresses these concerns; the negative controls of
  Section 10.5 provide a direct diagnostic.

  13.3 Absence of Physical Derivation
  The framework provides no dynamical mechanism, symmetry principle, or
  field-theoretic foundation for the observed compressibility. All results are interpretable
  only as empirical compressibility under constrained symbolic grammars.

  13.4 Finite Dataset Limitations
  The target set T contains N < 15 constants (after excluding those with relative
  uncertainty > 10-2). Statistical conclusions from such small N are subject to substantial
  finite-sample variance. Asymptotic coverage theory applies only approximately.

  13.5 Identifiability and Non-Uniqueness of Fits
  Multiple distinct symbolic expressions may achieve comparable accuracy: F1 approx X
  approx F2 with F1 $\neq$ F2. This structural non-uniqueness is a property of all
  over-complete dictionaries and implies that the best-fit expression F\textasciicircum{}*(X) is not
  uniquely defined near dense regions of the grammar lattice.

  13.6 Asymptotic Saturation
  As C $\rightarrow$ $\in$fty, S(C) becomes dense in R\textasciicircum{}+ and coverage trivially approaches 1 for any $\varepsilon$
  > 0. The analysis is restricted to C $\leq$ 10 to avoid this uninformative regime.

  13.7 Sensitivity to Complexity Metric Choice
  The ell1-norm complexity metric induces a hypothesis class that differs from those
  induced by ell0, ell2, or depth-based metrics. All reported results are specific to the ell1
  choice. Sensitivity analyses with alternative metrics are reserved for future work.

  13.8 Experimental Uncertainty Propagation
  For constants with relative uncertainty sigmaXi/Xi > 10-2, a symbolic approximation at
  tolerance $\varepsilon$ = 10-3 is indistinguishable from the best achievable match given
  experimental precision. Such constants are flagged and excluded from the primary
  analysis.

